\section{Controller and wire-level details}

\subsection{Controller state machine}
Fig.~\ref{fig:fsm} shows the controller of \texttt{cgra\_ctrl} with the RTL
state names. From \texttt{S\_IDLE} the first received byte selects a branch.
Short transactions (\cmd{ID}, \cmd{RST} and the rejection of an unknown
opcode) reply at once through the shared transmit handshake. Bulk transfers
accumulate their payload, receive and verify the checksum in a dedicated
state and reply \cmd{ACK} or \cmd{NACK}. \cmd{RUN} produces $N$ enabled PE
captures at the configured cadence before acknowledging. The \cmd{EXEC}
header state falls through into the same receive state as \cmd{WR} and the
same step engine as \cmd{RUN}; only the return targets, held in two
registers, differ. A watchdog guards every state that waits for a payload
or checksum byte. After an uncertain transfer the library clears its cached
identity and fails later calls with \texttt{CGRA\_ERR\_DESYNC}; it sends no
implicit \cmd{RST}, because a controller waiting for a \cmd{RUN} argument
would interpret that byte (\byte{06}) as a step count.

\begin{figure*}[t]
  \centering
  \resizebox{\textwidth}{!}{% Controller state machine (hw/rtl/cgra_ctrl.vhd). Each box is one RTL state
% or a short chain of them (named in typewriter); byte loops and the shared
% transmit handshake S_SEND/S_SEND_WAIT are folded into the boxes that use it.
% Blue boxes are the states shared by two commands: S_WR (WR, EXEC) and
% S_RUN (RUN, EXEC); the return target is held in wr_ret / run_ret.
\begin{tikzpicture}[x=1mm, y=1mm, font=\scriptsize,
    st/.style={draw=cStateE, fill=cState, rounded corners=2pt,
      minimum width=27mm, minimum height=9.5mm, align=center, inner sep=1pt},
    sh/.style={st, draw=cBlockE, fill=cBlock},
    ex/.style={st, draw=cRegE, fill=cReg},
    tx/.style={st, draw=cWire!70, fill=white},
    sn/.style={font=\tiny\ttfamily, text=cWire!80},
    op/.style={lbl, fill=white, inner sep=1pt, font=\tiny\ttfamily}]
  \newcommand{\stbox}[5]{\node[#1] (#2) at #3 {#4\\[-0.4ex]{\tiny\ttfamily\color{cWire!80}#5}};}

  % hub
  \node[draw=cStateE, fill=cState, circle, minimum size=13mm, align=center,
        font=\scriptsize\bfseries] (idle) at (0,0) {IDLE\\[-0.5ex]{\tiny\ttfamily\mdseries S\_IDLE}};

  % one row per command family: y positions
  \def\yCFG{33} \def\yWR{20} \def\yRUN{7} \def\yRD{-6} \def\yID{-19} \def\yEX{-33}

  \stbox{st}{cfg1}{(36,\yCFG)}{receive $RC$ words}{S\_CFG}
  \stbox{st}{cfg2}{(70,\yCFG)}{receive, verify checksum}{S\_CFG\_CK}
  \stbox{tx}{cfg3}{(104,\yCFG)}{send ACK / NACK}{S\_SEND}

  \stbox{sh}{wr1}{(36,\yWR)}{receive $R{+}C$ words}{S\_WR}
  \stbox{st}{wr2}{(70,\yWR)}{receive, verify checksum}{S\_WR\_CK}
  \stbox{tx}{wr3}{(104,\yWR)}{send ACK / NACK}{S\_SEND}

  \stbox{st}{run1}{(36,\yRUN)}{receive $N$}{S\_RUN\_ARG}
  \stbox{sh}{run2}{(70,\yRUN)}{$N$ enabled captures}{S\_RUN}
  \stbox{tx}{run3}{(104,\yRUN)}{send ACK}{S\_RUN\_ACK}

  \stbox{tx}{rd1}{(36,\yRD)}{send $RC$ registers}{S\_RD}
  \stbox{tx}{rd2}{(70,\yRD)}{send checksum}{S\_RD\_CK}

  \stbox{tx}{id1}{(36,\yID)}{direct reply}{S\_ID / S\_SEND}
  \node[lbl, align=left, anchor=west] (idl) at (51,\yID)
    {ID: signature, version, geometry\\RST: clear PE registers, ACK\\unknown opcode: NACK};

  \stbox{ex}{ex1}{(36,\yEX)}{header: $N$, flags, mask}{S\_EX\_HDR}
  \stbox{sh}{ex2}{(70,\yEX)}{receive $R{+}C$ words}{S\_WR}
  \stbox{ex}{ex3}{(104,\yEX)}{receive, verify checksum}{S\_EX\_CK}
  \stbox{sh}{ex4}{(138,\yEX)}{optional reset,\\[-0.4ex] $N$ captures}{S\_RUN}
  \stbox{ex}{ex5}{(172,\yEX)}{masked registers,\\[-0.4ex] checksum, status}{S\_EX\_RD\dots S\_EX\_ST}

  % dispatch from IDLE: orthogonal fan-out with opcode labels
  \foreach \y/\dst/\code in {\yCFG/cfg1/02\,CFG, \yWR/wr1/03\,WR, \yRUN/run1/04\,RUN,
                             \yRD/rd1/05\,RD, \yID/id1/{01,06,other}, \yEX/ex1/07\,EXEC} {
    \draw[wire] (idle.east) -| (14,\y) -- node[op, above=0.3mm] {\code} (\dst.west);
  }

  % forward edges
  \foreach \a/\b in {cfg1/cfg2, cfg2/cfg3, wr1/wr2, wr2/wr3, run1/run2, run2/run3,
                     rd1/rd2, ex1/ex2, ex2/ex3} \draw[wire] (\a) -- (\b);
  \draw[wire] (ex3) -- node[op, above=0.3mm] {valid} (ex4);
  \draw[wire] (ex4) -- (ex5);
  \draw[wire, cRegE] (ex3.south) -- ++(0,-6) -| node[op, pos=0.25, below=0.3mm, text=cRegE]
    {invalid: no reset, no steps} (ex5.south);

  % return to IDLE through one bus on the right and below
  \def\xb{192}
  \foreach \n in {cfg3, wr3, run3, rd2, idl, ex5} \draw[cWire, thick] (\n.east) -- (\xb, |- \n.east);
  \draw[cWire, thick] (\xb,\yCFG) -- (\xb,-47);
  \draw[wire] (\xb,-47) -- node[op, below=0.3mm] {reply complete} (0,-47) -- (idle.south);
  \draw[cWire, thick] (id1.east) -- (idl.west);

  % watchdog (dashed): from every receive state back to IDLE
  \draw[wire, dashed, cStateE] (128,44) -- node[op, above=0.3mm, text=cStateE]
    {watchdog: no byte for $10^{8}$ cycles (1\,s) in any state awaiting a payload or checksum byte}
    (0,44) -- (idle.north);
\end{tikzpicture}
}
  \caption{Controller state machine of \texttt{cgra\_ctrl}. Boxes give the
  function and the RTL state(s) in typewriter; byte loops and the shared
  transmit handshake are folded into the boxes. Blue states are shared by two
  commands (\cmd{WR}/\cmd{EXEC} reception, \cmd{RUN}/\cmd{EXEC} stepping) and
  return through a stored target; green states exist only for \cmd{EXEC};
  white states transmit. A rejected \cmd{EXEC} checksum skips reset and
  stepping but still returns the current registers with \cmd{NACK}. Every
  reply ends in \texttt{S\_IDLE}; the watchdog (dashed) aborts stalled
  reception.}
  \label{fig:fsm}
\end{figure*}


\subsection{A transaction on the wire}
Writing $\text{west}=(1,0,0,0)$ and $\text{north}=(5,0,0,0)$ on a $4\times4$
device sends the command byte \byte{03}, the four west and four north words
little-endian, and the additive checksum:
\begin{center}\small\ttfamily
03\quad 0100 0000 0000 0000\quad 0500 0000 0000 0000\quad 06
\end{center}
The sixteen payload bytes sum to 6; the device replies \cmd{ACK}
(\byte{79}) or, on a mismatch, \cmd{NACK} (\byte{1F}). Words are committed
as they arrive, so a \cmd{NACK} does not restore the previous inputs.
